% 智能体大语言模型的持续强化学习 —— 中文版
% 编译: xelatex main_cn; bibtex main_cn; xelatex x2

\documentclass[11pt]{article}

\usepackage{xeCJK}
\usepackage{amsmath,amssymb}
\usepackage{graphicx}
\usepackage{booktabs}
\usepackage{hyperref}
\usepackage{xcolor}

\graphicspath{{../assets/}}

\newcommand{\TODO}[1]{{\color{red}[TODO: #1]}}

\title{智能体大语言模型的持续强化学习\\
       \large 基于能力分区的经验回放与抗遗忘评测}

\author{孙豪}
\date{}

\begin{document}
\maketitle

\begin{abstract}
本文研究智能体大语言模型在按能力域顺序训练的持续强化学习问题——线上数据更新快、分布不均衡，无法一次集齐全部分布，只能逐个能力域顺序训练，这恰好放大了灾难性遗忘。为此本文设计了一套基于能力分区的九桶经验回放缓冲区：按能力域而非难度分桶，容量由任务数平方根加权分配，桶内淘汰禁止跨桶挤出，桶间按能力空间距离采样回放以确保遗忘风险越高的桶获得越多回放权重。持续学习目标由策略梯度、逆向KL锚定、回放损失和固定熵正则四项组成，以零框架改动注入GRPO。奖励由外部冻结大模型裁判统一打分，采用多维乘性聚合并以安全作为二元门控。评测方面，本文提出按桶分组、一次统一评测、权重后置的抗遗忘评测协议，量化"越早训练的桶遗忘越重"并与同规模无持续学习机制的基线对照。本文在Qwen3.5-9B上跨九个能力域实现并验证整套系统。
\end{abstract}

\section{引言}
\label{sec:intro}

线上部署的大语言模型智能体通过周期性强化学习更新持续进化：每隔一段时间，用新采集的交互数据精炼策略。这本质上是\textbf{持续强化学习}问题——模型必须在习得新能力的同时保住已有能力。与实验室一次性集齐全分布数据不同，线上数据更新快、到达不均衡，往往只能\textbf{按能力域分批、顺序地}训练。这恰好放大了\textbf{灾难性遗忘}~\cite{D1}：后训练的域会系统性侵蚀先训练的域已掌握的能力。

针对遗忘的主流解法是\textbf{经验回放}~\cite{A2}：训练新数据时从历史经验采样一部分回放进当前批次，用独立的持续学习损失让模型温习旧能力。已有工作的隐含假设是被回放的经验携带忠实的学习信号。在多能力域场景下，一个额外挑战浮现：不同能力会在扁平缓冲区中互相挤占，"回放什么、回放多少"本身就是一个设计问题。

本文面向智能体大语言模型的持续强化学习，设计并实现一套完整的训练与评测系统，核心贡献有三：

\begin{enumerate}
  \item \textbf{按能力分区的九桶经验回放缓冲区}。桶按能力域而非难度划分，每桶设固定容量上限（任务数平方根加权）和保护下限。桶内淘汰、禁止跨桶挤出，桶间按能力空间\textbf{距离}采样回放——离当前训练批越远的桶遗忘风险越高、回放权重越大。桶内采样和淘汰策略为可替换设计，当前分别采用随机采样和FIFO淘汰。

  \item \textbf{面向持续学习的复合训练目标}，以零框架改动注入GRPO：策略梯度、对参考策略的逆向KL锚定、经验回放损失、固定开启的熵正则（防策略多样性坍缩~\cite{B4}）。奖励由单个外部冻结大模型裁判在统一细则下打分，多维乘性聚合，安全作为二元门控以提高评价一致性。

  \item \textbf{按桶分组的抗遗忘评测协议}：全部能力域顺序训完后一次统一评测，施加按训练先后衰减的权重后置量化"越早训练的桶遗忘越重"，并与同规模无持续学习机制的基线对照。
\end{enumerate}

本文在Qwen3.5-9B智能体上跨九个能力域实现并验证整套系统，用Pass$^3$基准~\cite{C1}同时刻画新能力习得与旧能力遗忘。

\section{方法}
\label{sec:method}

\subsection{持续学习复合目标}

在标准GRPO策略梯度之外叠加三项持续学习正则，构成复合损失：

\begin{equation}
  L_{cl} = \lambda_1 L_{rl} + \lambda_2 L_{kl} + \lambda_3 L_{replay} + \lambda_4 L_{ent}.
  \label{eq:cl-loss}
\end{equation}
$L_{rl}$为GRPO项（组内奖励归一化估计优势）。$L_{kl}=D_{KL}(\pi_{new}\,\|\,\pi_{ref})$为对参考策略$\pi_{ref}$的逆向KL锚定。$L_{replay}$为经验回放损失，作用于从缓冲区采样的历史轨迹（第\ref{sec:buffer}节）；回放行携带独立的\texttt{replay\_response\_mask}，不参与PPO主损失。$L_{ent}$为熵正则，$\lambda_4=0.001$所有实验固定开启，防策略多样性坍缩~\cite{B4}。早期参数L2正则$L_{reg}$已弃用。目标以\textbf{零框架改动}注入——通过框架的损失注入接口与钩子系统挂载。

\subsection{奖励设计}
\label{sec:reward}

由单个外部冻结大模型裁判在统一细则下打分，四个维度：

\begin{itemize}
  \item \textbf{Done}（完成度，$0/1$）：任务是否真正完成。
  \item \textbf{Correctness}（正确性，$[0,1]$）：产出是否正确。
  \item \textbf{Trajectory}（轨迹质量，$[0,1]$）：工具调用质量、有无冗余步、推理是否连贯；\textbf{并惩罚"加了不该加的东西"}。
  \item \textbf{Safety}（安全，\textbf{二元$0/1$}）：是否出现破坏性/越权操作。二元门控而非连续打分，一致性更高。
\end{itemize}

最终奖励乘性聚合：

\begin{equation}
  r =
  \begin{cases}
    (0.4\cdot\text{correctness} + 0.4\cdot\text{trajectory} + 0.2)\cdot\text{safety}, & \text{done}=1,\\[4pt]
    0.4\cdot\text{trajectory}\cdot\text{safety}, & \text{done}=0.
  \end{cases}
  \label{eq:reward}
\end{equation}

安全作为$\{0,1\}$门控，一旦触发整条清零。未完成任务仍给轨迹分，但无完成基线分和正确性分。

\subsection{九桶经验回放缓冲区}
\label{sec:buffer}

\paragraph{按能力分桶。}桶按能力/领域划分。九个纯文本能力域（表\ref{tab:buckets}）由评测基准官方类目合并、去多模态得到。$K=9$桶，共195个纯文本任务，单层无子桶。

\begin{table}[t]
\centering
\caption{各桶任务数与容量上限。}
\label{tab:buckets}
\begin{tabular}{lrr}
\toprule
桶（能力域） & 任务数 & 容量上限 \\
\midrule
workflow（工作流编排） & 56 & 4915 \\
ops（系统操作） & 44 & 4354 \\
qa（问答检索） & 36 & 3938 \\
finance（财务金融） & 20 & 2935 \\
office（办公文档） & 11 & 2177 \\
communication（沟通表达） & 11 & 2177 \\
safety（安全合规） & 9 & 1969 \\
research（研究综合） & 6 & 1607 \\
coding（代码） & 2 & 928 \\
\midrule
合计 & 195 & 25000 \\
\bottomrule
\end{tabular}
\end{table}

\paragraph{容量分配。}各桶容量$\text{cap}_i$由任务数$n_i$的平方根加权分配，总容量$C=25000$：

\begin{equation}
  \text{cap}_i = C\cdot\frac{n_i^{\alpha}}{\sum_j n_j^{\alpha}},\qquad \alpha=0.5.
  \label{eq:quota}
\end{equation}
取$\alpha=0.5$（次线性）使大桶温和倾斜、小桶被抬起。每桶另设保护下限（约为cap的$1/30$）。

\paragraph{淘汰与采样。}淘汰\textbf{仅在桶内进行}——禁止跨桶挤出。默认淘汰策略为FIFO。CLEAR基线用蓄水池采样~\cite{A2}。回放采样\textbf{两级}：先按能力空间\textbf{距离}选桶，离当前训练批质心越远权重越大（遗忘风险越高）；桶内随机（uniform）采样。每步回放行总数固定。

\paragraph{冷启动种子。}预采集的完整轨迹（含工具调用与思考）可预填缓冲区。

\subsection{回放损失加权}
\label{sec:replay-weight}

每条回放轨迹在$L_{replay}$中按token级权重：

\begin{equation}
  w_t^{(i)} = \text{normalize}\!\Big(\text{clip}\big(\text{priority}_i \cdot \tfrac{\gamma^{\text{block}(t)} + \delta^{K_i - \text{block}(t)}}{2},\; q_5,\; q_{95}\big)\Big),
  \label{eq:u-weight}
\end{equation}
其中块按OpenAI chat message边界划分。第二个因子为U形块权重的一般形式。

\textbf{当前默认关闭U形：}$\gamma=\delta=1$，块权重退化为常数1.0，退化为纯priority加权（scheme W0）。U形（$\gamma=\delta<1$，scheme W2）作为独立消融项保留。

\section{实验}
\label{sec:experiments}

\subsection{实验设置}

\textbf{评测基准。}采用ClawEval的纯文本子集，共195个任务，按表\ref{tab:buckets}的九个能力域分组。评分采用Pass$^3$协议（三次独立运行全部通过），逐任务二元判定、逐桶通过率。

\textbf{训练协议。}训练\textbf{按桶顺序}进行：模型先在一个能力域上训练，再切换到下一个，模拟线上无法一次集齐全分布的现实约束。同一桶内任务随机打乱。基座模型为Qwen3.5-9B，采用完全异步的策略分离部署（推理与训练分卡），BF16全栈。

\textbf{算法配置。}策略采用GRPO，每组8条查询轨迹。裁判模型为独立部署的冻结外部LLM。底层RL框架verl 0.8.0不做任何改动，所有持续学习逻辑通过损失注入接口与缓冲区钩子挂载。

\subsection{抗遗忘评测协议}

全部能力域顺序训完后\textbf{一次统一评测}（训练过程中不做中间评测，省算力）：

\begin{enumerate}
  \item 对终态策略$\pi_K$跑一次全量评测。
  \item 评测结果按同样的九个桶分组，组间序与训练序对齐，记录每桶原始分数$S_i = S(\pi_K, b_i)$。
  \item 施加按训练先后\textbf{衰减的权重}$w_i$：越早训练的桶遗忘越重、权重越大：
  \begin{equation}
    \text{CL-Score} = \sum_{i=1}^{K} w_i\,S_i,\qquad w_i \propto g(i),\;\;g\text{递减}.
  \end{equation}
  衰减形式（线性/指数）作为超参由实验确定。
\end{enumerate}

评测只跑一次，之后切换权重方案只是算术运算，无需重新推理。核心信号是含持续学习机制的模型在\textbf{更早训练的桶上}相对基线表现出\textbf{更小的分数下降}。

\subsection{对照与消融}

\textbf{核心对照}：含持续学习机制（九桶回放+复合损失）\emph{vs.} 无持续学习机制的同规模基线（无回放、$\lambda_3=0$，相同训练序和评测）。

\textbf{消融轴}：
\begin{itemize}
  \item 缓冲区结构：九桶分区回放 \emph{vs.} CLEAR式单扁平缓冲区+蓄水池采样~\cite{A2}。
  \item 缓冲区容量：不同总容量下的抗遗忘表现。
  \item U形加权：$\gamma=\delta=1$（均权，默认）\emph{vs.} $\gamma=\delta<1$（U形激活）。
\end{itemize}

\subsection{实验结果}

\todo{新reward+CL算法16卡正式实验就绪后填入。预期产出：逐桶原始分数；CL-Score对照表（CL vs. baseline）；reward/loss/组内std训练曲线；消融柱状图。}


\bibliographystyle{plain}
\bibliography{../refs/references}

\appendix
\section{奖励裁判打分细则}
\label{app:prompts}

奖励裁判（冻结外部LLM）每条轨迹接收三个输入：任务描述、智能体轨迹、打分细则。系统提示固定输出格式：

\begin{quote}\small
You are a strict evaluator for autonomous-agent trajectories. Think step by step, then grade the agent on the four dimensions defined in the rubric.

Output ONLY a JSON object with keys task\_done, correctness, trajectory, safety. task\_done is 0 or 1; safety is 0 or 1; correctness and trajectory are floats in [0,1]. No prose, no explanation, no markdown code fences.

Example: \{"task\_done": 1, "correctness": 0.5, "trajectory": 0.8, "safety": 1\}.
\end{quote}

\subsection*{task\_done（完成度，0/1）}
智能体是否真正完成了任务？产出文件类任务以环境差异中的真实产物为准（非声称）；无文件产出时从最终回答判断。

\subsection*{correctness（正确性，$[0,1]$）}
产出是否正确？提供答案键时对照已知正确值——匹配\textbf{值}而非顺序，容忍小幅四舍五入。无答案键时由裁判据实判断。

\subsection*{trajectory（轨迹质量，$[0,1]$）}
三个子维度：（1）工具调用质量——格式正确、成功执行的调用得分；（2）无冗余步——惩罚无意义重复，\textbf{并惩罚"加了不该加的东西"}；（3）推理连贯。

\subsection*{safety（安全，二元0/1）}
\textbf{二元门控。}是否出现了有害/越权/破坏性操作？边界但无害的行为归入轨迹质量评定。1=安全，0=发生了破坏性操作。

\subsection*{聚合}
提取JSON判词后按下式计算最终奖励：

\begin{equation}
  r =
  \begin{cases}
    (0.4\cdot c + 0.4\cdot t + 0.2)\cdot s, & \text{done}=1,\\[4pt]
    0.4\cdot t\cdot s, & \text{done}=0,
  \end{cases}
\end{equation}
其中$c=$ correctness，$t=$ trajectory，$s=$ safety $\in\{0,1\}$。JSON解析失败重试一次，再失败则丢弃该轨迹。


\end{document}
